% !PRINTED 114
\weberpage{114}
%<<CONT

\begin{equation}
 \varepsilon_k=\cos\frac{2k\pi}{n}+i\sin\frac{2k\pi}{n}.
\tag{11}
\end{equation}
The quantities $\varepsilon_k$ are roots of the equation
\begin{equation}
 x^n=1,
\tag{12}
\end{equation}
and are called the $n$th \gesperrt{roots of unity}. There are $n$, and only $n$, different values of them; for odd $n$ only one is real, namely for $k=0$, and for even $n$ two are real, namely for $k=0$ and $k=\frac n2$. The geometrical images of the $n$th roots of unity are the vertices of a regular $n$-gon inscribed in the circle of radius $1$. Their algebraic determination is the subject of the \gesperrt{theory of cyclotomy}. If we denote by $\varepsilon$ the value of $\varepsilon_k$ for $k=1$, thus put
\[
 \varepsilon=\cos\frac{2\pi}{n}+i\sin\frac{2\pi}{n},
\]
then, by \gesperrt{Moivre's} theorem,
\[
 \varepsilon_k=\varepsilon^k,
\]
so that all $n$th roots of unity are represented as powers of one among them.


\begin{center}
{\large \S.\ 34.}

\vspace{0.5em}
{\large\gesperrt{Removal of the Second Term of an Equation.}}
\end{center}

The certainty of the existence of the roots of pure equations puts us in a position to prove the existence of roots for equations of the second, third, and fourth degrees, or, as one also says, for quadratic, cubic, and biquadratic equations, by reducing the determination of their roots to the solution of pure equations. We shall return to this question later from various more general points of view, but here we shall briefly discuss the older methods of solution, which, even if they give little insight into the general connection of this question, are nevertheless very simple in application. They arouse the appearance as if one were dealing with a method capable of generalization to higher equations, which is not the case. First the following general remark.

Let us, for the sake of simplicity, take in the function
%CONT>>
